\documentclass[11pt]{article}
\usepackage[a4paper,margin=26mm]{geometry}
\usepackage{amsmath,amssymb,mathtools}
\usepackage{unicode-math}
\setmainfont{Latin Modern Roman}
\setmathfont{Latin Modern Math}
\usepackage{xurl,hyperref}
\hypersetup{hidelinks,pdftitle={Problem statement}}
\newcommand{\sref}[2]{\hyperlink{source:#1:#2}{[S#2]}}
\newcommand{\cataloguescope}[1]{#1}
\begin{document}
\section{Foliated Teichmüller theory — Question 11.1}\label{statement}

\subsection{Definitions and mathematical statement}

\paragraph{Shared notation.}$\mathbb N=\{1,2,\ldots\}$, and $\mathbb Z,\mathbb Q,\mathbb R,\mathbb C$ denote the integers, rationals, reals and complex numbers. Any other conventions are specified locally.


A codimension-one foliation $\mathcal F$ decomposes a $3$-manifold $M$ into surface leaves, locally as the horizontal slices of a product chart. In the terminology of the source, it is taut if a closed curve transverse to $\mathcal F$ meets every leaf. A minimal foliation means that every leaf is dense; a minimal surface means a surface with vanishing mean curvature for the specified Riemannian metric. These two uses of minimal are distinct. A manifold is atoroidal if it has no incompressible embedded torus other than a boundary-parallel one. Write $\widetilde M$ for its universal cover, and put a tilde on a foliation or lamination for its lift. Consider automorphisms of a foliated pair $(M,\mathcal F)$ taking leaves to leaves, and their classes modulo isotopies through such automorphisms. This is the foliation mapping-class viewpoint; it can be richer than the ordinary mapping class group of $M$.
\paragraph{Statement recorded in the supplied source.}
Determine what nontrivial mapping-class elements can occur for taut foliations, including their geometric realization and possible infinite-order behavior. In the slithering example of the source a transverse flow can have a time-one map preserving the foliation, even though intermediate times need not preserve leaves; such elements must not be discarded merely because the underlying ambient map is isotopic to the identity. \sref{statement}{2}
\subsection{Short English statement}
Classify nontrivial foliation-preserving mapping classes, including those produced by a transverse flow and invisible to the ordinary ambient mapping class group.
\subsection{Sources}
\begin{description}
\item[\hypertarget{source:statement:1}{S1}] ULAM.AI and the UnsolvedMath Contributors, \emph{UnsolvedMath: A Curated Collection of Open Mathematics Problems}, catalogue entry. CC BY 4.0. \url{https://huggingface.co/datasets/ulamai/UnsolvedMath}.
\item[\hypertarget{source:statement:2}{S2}] Danny Calegari, \emph{Problems in foliations and laminations of 3-manifolds}, version 0.78 (2002), Question 11.1; Definitions 1.1--1.10 and local definitions and remarks. \url{https://arxiv.org/abs/math/0209081v1}.
\end{description}
\label{end:statement}
\end{document}
